\section{Conclusion and Future Work}
\label{sec:conclusion}

This study investigated whether a shared multi-dataset selective-state-space encoder for rotating-machinery fault diagnosis can be compressed without an unacceptable loss of diagnostic performance. Rather than designing an unrelated lightweight classifier, the proposed K1 model is structurally derived from the bidirectional Full-S1 PC-STE encoder by removing the reverse selective-scan path while retaining four temporal stages, the 192-dimensional representation, physical-coordinate encoding, cross-band mixing, and dataset-specific output heads. K1 is then distilled from a same-fold ensemble of independently seeded Full-S1 teachers.

The sealed nine-cell evaluation shows that this structural reduction preserves the diagnostic role of the shared encoder. Macro-4 Macro-F1 changes from $0.934644\pm0.023258$ for Full-S1 to $0.955334\pm0.018393$ for K1, with a mean paired difference of $+0.020691\pm0.013509$. The prespecified $-0.02$ Macro-F1 non-inferiority criterion is satisfied by the exact paired sign-flip test. At the same time, K1 reduces encoder parameters by $42.24\%$, FP32 serialized size by $42.14\%$, and equal-domain computation by $45.38\%$. Under the corrected inference benchmark, K1 provides approximately $1.87\times$ CPU and $1.63\times$ GPU end-to-end batch-1 speed-ups, and nearly doubles GPU batch-32 throughput.

The secondary Q8 extension further separates storage compression from architectural acceleration. The registered simulated-Q8 representation retains essentially the same Macro-4 Macro-F1 as K1 and satisfies its historical $-0.01$ non-inferiority criterion, while reducing stored model size by $71.12\%$ relative to K1 FP32 and by $83.29\%$ relative to Full-S1. The corresponding real CPU dynamic-quantization speed-up is modest because the selective-scan path and other sensitive components remain FP32. Q8 should therefore be interpreted primarily as an additional storage-compression mechanism on the current stack, not as evidence of GPU INT8 acceleration.

The broader implication is that deployment efficiency can be improved by compressing the structure of a reusable representation rather than replacing it with a separate task-specific small model. In this experiment, directionality-specific reduction accounts for most of the runtime benefit, while post-training quantization provides additional storage savings. The matched fold--seed design and prespecified non-inferiority criterion provide an explicit way to judge whether those savings are obtained at an acceptable predictive cost.

Future work should evaluate the same models with fused or hardware-optimized selective-scan kernels and on physical edge devices, where the relationship between theoretical FLOPs and realized latency may differ from the present workstation benchmark. The exact CPU dynamic-Q8 artifact should also receive its own pre-registered predictive evaluation if a deployment-level accuracy claim is required. More importantly, the shared encoder should be tested under genuinely unseen-machine, unseen-rig, and unseen-acquisition-system protocols, rather than only within-dataset held-out partitions. Finally, future representation-learning work can target the remaining class-separation weaknesses, particularly on CWRU, and examine whether broader multi-dataset pretraining or multimodal vibration--language objectives can improve transfer while preserving the compact K1 deployment profile.
